% ============================================================
%  يُترجَم بـ XeLaTeX  (في Overleaf: Menu → Compiler → XeLaTeX)
%  الخط: Amiri (متوفر في Overleaf، أو حمّله من github.com/aliftype/amiri)
% ============================================================
\documentclass[12pt,a4paper]{article}
\usepackage[a4paper,margin=1.7cm,top=2.3cm,bottom=2.9cm,headheight=15pt,footskip=1.5cm]{geometry}
\usepackage{amsmath,amssymb}
\usepackage{icomma}
\usepackage[english,bidi=default,layout=lists]{babel}
\babelprovide[import,main]{arabic}
\babelfont{rm}[Scale=1.08]{Amiri}
\babelfont{sf}[Scale=1.08]{Amiri}
\usepackage[most]{tcolorbox}
\usepackage{tikz}
\usetikzlibrary{arrows.meta,calc,shadings}
\usepackage{array}
\usepackage{colortbl}
\usepackage{fancyhdr}
\usepackage{titlesec}
\usepackage{enumitem}
\usepackage{collcell}
\usepackage{needspace}
\usepackage{etoolbox}
\newcolumntype{C}{>{\collectcell\ar}c<{\endcollectcell}}
\pretocmd{\section}{\Needspace{15\baselineskip}}{}{}
\pretocmd{\subsection}{\Needspace{11\baselineskip}}{}{}

\definecolor{cDef}{HTML}{1F4E9C}
\definecolor{cThm}{HTML}{B3261E}
\definecolor{cRem}{HTML}{2E7D32}
\definecolor{cEx}{HTML}{C76B00}
\definecolor{cMeth}{HTML}{6A1B9A}
\definecolor{cSol}{HTML}{00695C}
\definecolor{cSec}{HTML}{1F3A68}

\titleformat{\section}{\Large\bfseries\color{cSec}}{\thesection.}{0.6em}{}[{\color{cSec!60}\titlerule[1pt]}]
\titleformat{\subsection}{\large\bfseries\color{cThm!85!black}}{\thesubsection.}{0.5em}{}
\titlespacing*{\section}{0pt}{14pt}{8pt}
\titlespacing*{\subsection}{0pt}{10pt}{5pt}
\setlength{\parindent}{0pt}
\setlength{\parskip}{4pt}
\setlist{itemsep=2pt,topsep=3pt}

\newtcolorbox{arbox}[3][]{enhanced,
  colback=#2!5,colframe=#2!80!black,arc=2mm,boxrule=0.8pt,
  left=6pt,right=6pt,top=8pt,bottom=5pt,
  fonttitle=\bfseries,coltitle=white,
  title={\foreignlanguage{arabic}{#3}},
  attach boxed title to top right={xshift=-5mm,yshift=-3mm},
  boxed title style={colback=#2!80!black,arc=1.5mm,boxrule=0pt},
  before upper={\selectlanguage{arabic}},
  before skip=8pt,after skip=8pt,breakable,#1}
\newenvironment{BOX}[2]{\begin{otherlanguage}{english}\begin{arbox}{#1}{#2}}{\end{arbox}\end{otherlanguage}}
\newenvironment{tarif}[1][تعريف]{\begin{BOX}{cDef}{#1}}{\end{BOX}}
\newenvironment{khassia}[1][خاصية]{\begin{BOX}{cThm}{#1}}{\end{BOX}}
\newenvironment{molahada}[1][ملاحظة]{\begin{BOX}{cRem}{#1}}{\end{BOX}}
\newenvironment{mithal}[1][مثال]{\begin{BOX}{cEx}{#1}}{\end{BOX}}
\newenvironment{tariqa}[1][طريقة]{\begin{BOX}{cMeth}{#1}}{\end{BOX}}

\newcounter{exo}
\newenvironment{tamrin}[1]{\stepcounter{exo}\begin{BOX}{cSec}{التمرين \theexo: #1}}{\end{BOX}}
\newenvironment{hal}{\Needspace{12\baselineskip}\begin{otherlanguage}{english}\begin{arbox}[breakable]{cSol}{حل التمرين \theexo}}{\end{arbox}\end{otherlanguage}}

\newenvironment{LTR}{\begin{otherlanguage}{english}}{\end{otherlanguage}}
\newcommand{\ar}[1]{\foreignlanguage{arabic}{#1}}

\renewcommand{\dots}{\mathinner{\ldotp\ldotp\ldotp}}
\renewcommand{\ldots}{\mathinner{\ldotp\ldotp\ldotp}}
\renewcommand{\cdots}{\mathinner{\cdotp\cdotp\cdotp}}
\newcommand{\N}{\mathbb{N}}
\newcommand{\Z}{\mathbb{Z}}
\newcommand{\D}{\mathbb{D}}
\newcommand{\Q}{\mathbb{Q}}
\newcommand{\R}{\mathbb{R}}
\newcommand{\pgcd}{\operatorname{PGCD}}
\newcommand{\ppcm}{\operatorname{PPCM}}

% \NLseg{a}{b}{[}{]} ; \NLright{a}{[} ; \NLleft{b}{]}  (interval on a number line)
\newcommand{\NLseg}[4]{\begin{LTR}\begin{tikzpicture}[baseline=-0.6ex]
 \draw[-{Stealth[length=2mm]},gray!70](0,0)--(3.4,0);
 \draw[line width=2pt,cThm!80!black](0.7,0)--(2.7,0);
 \node[cThm!80!black,font=\LARGE] at (0.7,0){$#3$};
 \node[cThm!80!black,font=\LARGE] at (2.7,0){$#4$};
 \node[below=3pt,font=\footnotesize] at (0.7,0){$#1$};
 \node[below=3pt,font=\footnotesize] at (2.7,0){$#2$};
\end{tikzpicture}\end{LTR}}
\newcommand{\NLright}[2]{\begin{LTR}\begin{tikzpicture}[baseline=-0.6ex]
 \draw[-{Stealth[length=2mm]},gray!70](0,0)--(3.4,0);
 \draw[line width=2pt,cThm!80!black](1.3,0)--(3.3,0);
 \node[cThm!80!black,font=\LARGE] at (1.3,0){$#2$};
 \node[below=3pt,font=\footnotesize] at (1.3,0){$#1$};
\end{tikzpicture}\end{LTR}}
\newcommand{\NLleft}[2]{\begin{LTR}\begin{tikzpicture}[baseline=-0.6ex]
 \draw[-{Stealth[length=2mm]},gray!70](0,0)--(3.4,0);
 \draw[line width=2pt,cThm!80!black](0.05,0)--(2.1,0);
 \node[cThm!80!black,font=\LARGE] at (2.1,0){$#2$};
 \node[below=3pt,font=\footnotesize] at (2.1,0){$#1$};
\end{tikzpicture}\end{LTR}}

% ---------- social-media icons (pure TikZ) ----------
\definecolor{cFB}{HTML}{1877F2}
\definecolor{cTG}{HTML}{229ED9}
\definecolor{cIG1}{HTML}{FEDA75}
\definecolor{cIG2}{HTML}{D62976}
\definecolor{cIG3}{HTML}{4F5BD5}
\newcommand{\fbicon}{\begin{LTR}\begin{tikzpicture}[baseline=-0.75ex,x=1em,y=1em]
 \fill[cFB] (0,0) circle (0.55);
 \node[white,inner sep=0,font=\usefont{TU}{lmss}{bx}{n}\small] at (0.04,-0.02) {f};
\end{tikzpicture}\end{LTR}}
\newcommand{\igicon}{\begin{LTR}\begin{tikzpicture}[baseline=-0.75ex,x=1em,y=1em]
 \shade[lower left=cIG1,lower right=cIG2,upper left=cIG2,upper right=cIG3,rounded corners=0.28em] (-0.55,-0.55) rectangle (0.55,0.55);
 \draw[white,line width=0.09em,rounded corners=0.18em] (-0.36,-0.36) rectangle (0.36,0.36);
 \draw[white,line width=0.09em] (0,0) circle (0.17);
 \fill[white] (0.22,0.22) circle (0.05);
\end{tikzpicture}\end{LTR}}
\newcommand{\tgicon}{\begin{LTR}\begin{tikzpicture}[baseline=-0.75ex,x=1em,y=1em]
 \fill[cTG] (0,0) circle (0.55);
 \fill[white] (-0.33,0.0) -- (0.32,0.27) -- (0.2,-0.3) -- (0.0,-0.15) -- (-0.08,-0.28) -- (-0.08,-0.08) -- cycle;
 \fill[cTG!70!black] (-0.08,-0.08) -- (0.25,0.2) -- (0.0,-0.15) -- (-0.08,-0.28) -- cycle;
\end{tikzpicture}\end{LTR}}
% page number: black box, white two-digit number, at the right of the social row
\newcommand{\PageBox}{\begin{LTR}\tikz[baseline=(p.base)]\node[fill=black,text=white,inner sep=2.5pt,minimum width=1.6em,minimum height=1.3em,font=\usefont{TU}{lmss}{bx}{n}\small](p){\ifnum\value{page}<10 0\fi\arabic{page}};\end{LTR}}
% visual left-to-right order: Telegram | Instagram | Facebook | page box
\newcommand{\SocialRow}{\begin{otherlanguage}{english}%
 \ar{\textbf{الأستاذ سوايسية هشام للرياضيات}}\ \tgicon\hfill
 \textbf{prof\_hicham\_math}\ \igicon\hfill
 \ar{\textbf{الأستاذ سوايسية هشام للرياضيات}}\ \fbicon\hspace{2.2em}\PageBox\end{otherlanguage}}

\pagestyle{fancy}
\fancyhf{}
\renewcommand{\headrulewidth}{0.6pt}
\renewcommand{\footrulewidth}{0.4pt}
% \SetHeader{right head: document title}{left head: level -- stream}
\newcommand{\SetHeader}[2]{%
 \fancyhead[C]{\parbox{\textwidth}{\ar{#1\hfill #2}}}%
 \fancyfoot[C]{\parbox{\textwidth}{\small\SocialRow}}}
\usepackage{tkz-tab}
\usepackage{multicol}
% variation / sign table (LTR)
\newenvironment{TAB}{\par\begin{center}\begin{LTR}\begin{tikzpicture}[scale=0.95,every node/.style={font=\small}]}{\end{tikzpicture}\end{LTR}\end{center}}
% grid + axes, clipped: \begin{GR}{xmin}{xmax}{ymin}{ymax}{unit cm}
\newenvironment{GR}[5]{\par\begin{center}\begin{LTR}\begin{tikzpicture}[x=#5cm,y=#5cm,font=\footnotesize]
 \draw[gray!25,very thin] (#1,#3) grid (#2,#4);
 \draw[-{Stealth[length=2mm]}] (#1,0)--(#2+0.3,0) node[below left]{$x$};
 \draw[-{Stealth[length=2mm]}] (0,#3)--(0,#4+0.3) node[below left]{$y$};
 \node[below left] at (0,0){$O$};
 \begin{scope}\clip (#1,#3) rectangle (#2,#4);}
{\end{scope}\end{tikzpicture}\end{LTR}\end{center}}
\newcommand{\Ticks}[2]{\foreach \t in {#1} \node[below,font=\scriptsize] at (\t,0){$\t$}; \foreach \t in {#2} \node[left,font=\scriptsize] at (0,\t){$\t$};}
\newcommand{\bx}[1]{\ensuremath{\colorbox{cDef!12}{$\displaystyle #1$}}}
\newcommand{\TwoCol}[4]{\par\begin{LTR}\noindent\begin{minipage}[t]{#1\textwidth}\selectlanguage{arabic}#2\end{minipage}\hfill\begin{minipage}[t]{#3\textwidth}\selectlanguage{arabic}#4\end{minipage}\end{LTR}\par}
\newcommand{\RL}[4]{\TwoCol{#3}{#4}{#1}{#2}}  % \RL{w1}{right content}{w2}{left content}
\newcommand{\ok}{$\textcolor{cRem}{\checkmark}$}
\newcommand{\limn}{\displaystyle\lim_{n\to+\infty}}
% sub-questions, each on its own line
\newenvironment{sq}{\begin{itemize}[leftmargin=2.4em,itemsep=3pt]}{\end{itemize}}
\newcommand{\qa}{\item[(أ)]}\newcommand{\qb}{\item[(ب)]}\newcommand{\qc}{\item[(جـ)]}\newcommand{\qd}{\item[(د)]}\newcommand{\qe}{\item[(هـ)]}
\newcommand{\pa}{\par\smallskip\textbf{(أ)} }\newcommand{\pb}{\par\smallskip\textbf{(ب)} }\newcommand{\pc}{\par\smallskip\textbf{(جـ)} }\newcommand{\pd}{\par\smallskip\textbf{(د)} }\newcommand{\pe}{\par\smallskip\textbf{(هـ)} }
\newcommand{\QQ}[1]{\par\smallskip\textbf{#1.} }
\newenvironment{tmr}{\stepcounter{exo}\begin{BOX}{cSec}{تمرين \theexo}}{\end{BOX}}
\newcommand{\ABC}[3]{\\[2pt]\begin{tabular}{p{0.28\textwidth}p{0.28\textwidth}p{0.28\textwidth}}(جـ) \ #3 & (ب) \ #2 & (أ) \ #1\end{tabular}}
\pretocmd{\section}{\Needspace{24\baselineskip}}{}{}
\newtcbox{\bxtbox}{on line,colback=cDef!12,colframe=cDef!12,boxrule=0pt,arc=1pt,boxsep=0pt,left=2pt,right=2pt,top=1pt,bottom=1pt}
\newcommand{\bxt}[1]{\foreignlanguage{english}{\bxtbox{\ar{#1}}}}
\SetHeader{\small\bfseries\color{cSec}تمارين البكالوريا: الاحتمالات}{\small\color{cSec}السنة الثالثة ثانوي -- الشعب العلمية}
\begin{document}
\begin{LTR}
\begin{tcolorbox}[enhanced,colback=cSec,colframe=cSec,arc=3mm,halign=center,boxrule=0pt,top=12pt,bottom=12pt,before upper={\selectlanguage{arabic}}]
{\color{white}\fontsize{26}{32}\selectfont\bfseries تمارين البكالوريا في الاحتمالات}\\[8pt]
{\color{white}\large -- بدون حلول (للتدرّب) -- حتى دورة 2026}\\[5pt]
{\color{white}\large علوم تجريبية -- تقني رياضي -- رياضيات}
\end{tcolorbox}
\end{LTR}
\section{تمارين البكالوريا -- شعبة علوم تجريبية}
% ===== تمارين البكالوريا في الاحتمالات -- علوم تجريبية =====
\tikzset{
 bsT/.style={grow'=right,level distance=3.5cm,
   level 1/.style={sibling distance=2.8cm},level 2/.style={sibling distance=1.5cm},
   edge from parent/.style={draw,thick,cDef!70!black}},
 bsev/.style={circle,draw=cDef,fill=cDef!12,thick,minimum size=7mm,inner sep=1pt,font=\small},
 bslf/.style={rectangle,rounded corners=3pt,draw=cDef,fill=cDef!12,thick,minimum width=11mm,minimum height=7mm,inner sep=1pt,font=\small},
 bsp/.style={pos=0.68,fill=white,inner sep=0.8pt,font=\footnotesize},
 bspr/.style={anchor=west,font=\small,xshift=3mm},
 bshd/.style={anchor=west,font=\small\bfseries,xshift=3mm,yshift=8mm,text=cSec}}

%======================== 1 ========================
\begin{tamrin}{بكالوريا 2026 -- الموضوع الأول -- علوم تجريبية}
\textbf{I)} يحتوي كيس على $9$ كريات متماثلة لا نفرق بينها باللمس، منها: كريتان تحملان الرقم $1$، وأربع كريات تحمل الرقم $2$، وثلاث كريات تحمل الرقم $3$.

نسحب عشوائيًا من الكيس ثلاث كريات في آن واحد، ونعتبر الحوادث:

$A$: «سحب ثلاث كريات تحمل نفس الرقم»، \quad $B$: «سحب ثلاث كريات أرقامها مختلفة مثنى مثنى»،

$C$: «سحب ثلاث كريات مجموع أرقامها عدد زوجي».

\QQ{1}
\begin{sq}
\qa أحسب $P(A)$، $P(B)$ و$P(C)$ احتمالات الحوادث $A$، $B$ و$C$ على الترتيب.
\qb بيّن أن: $P(A\cap C)=\dfrac{1}{21}$ ثم استنتج $P_A(C)$.
\end{sq}
\QQ{2} $X$ المتغير العشوائي الذي يرفق بكل عملية سحب لثلاث كريات عدد الكريات التي تحمل الرقم $2$.

عيّن قانون احتمال المتغير العشوائي $X$ ثم أحسب أمله الرياضياتي.

\textbf{II)} نضيف إلى محتوى الكيس السابق $n$ كرية تحمل الرقم $1$ ثم نسحب منه كريتين على التوالي دون إرجاع.

عيّن قيمة $n$ حتى يكون احتمال الحصول على كريتين لا تحمل أي منهما الرقم $1$ يساوي $\dfrac{1}{5}$.
\end{tamrin}

%======================== 2 ========================
\begin{tamrin}{بكالوريا 2026 -- الموضوع الثاني -- علوم تجريبية}
\textbf{I)} يحتوي كيس على $10$ كريات متماثلة لا نفرق بينها باللمس، منها ذات الألوان: الأخضر، الأحمر والأبيض.
نسحب من الكيس بطريقة عشوائية كرية واحدة، ونعرف قانون احتمال هذه التجربة بالجدول الآتي، حيث $\alpha\in\R$:
\begin{center}
\begin{tabular}{|C|C|C|C|}\hline
\rowcolor{cSec!12} الكرية المسحوبة & خضراء & حمراء & بيضاء\\\hline
الاحتمال & $\dfrac{1}{10}$ & $\alpha$ & $2\alpha$\\[8pt]\hline
\end{tabular}
\end{center}
أحسب قيمة $\alpha$ ثم استنتج عدد كل من الكريات الخضراء، الحمراء، البيضاء.

\textbf{II)} نفرض أن الكيس يحتوي على كرية واحدة خضراء، $3$ كريات حمراء و$6$ كريات بيضاء.
نسحب عشوائيًا في آن واحد ثلاث كريات من هذا الكيس، ونعتبر الحوادث:

$A$: «الحصول على كرية حمراء واحدة فقط»، \quad $B$: «الحصول على كريتين بالضبط من نفس اللون»،

$C$: «الحصول على ثلاث كريات ليست كلها من نفس اللون».
\QQ{1}
\begin{sq}
\qa أحسب $P(A)$، $P(B)$ و$P(C)$ احتمالات الحوادث $A$، $B$ و$C$ على الترتيب.
\qb أحسب احتمال الحصول على كريتين بيضاوين علمًا أننا لا نحصل على الكرية الخضراء.
\end{sq}
\QQ{2} $X$ المتغير العشوائي الذي يرفق بكل عملية سحب لثلاث كريات عدد الألوان التي تحملها الكريات المسحوبة.

بيّن أن: $P(X=2)=\dfrac{27}{40}$، وعيّن قانون احتمال المتغير العشوائي $X$ ثم أحسب أمله الرياضياتي $E(X)$.
\end{tamrin}

%======================== 3 ========================
\begin{tamrin}{بكالوريا 2025 -- الموضوع الأول -- علوم تجريبية}
يحتوي صندوق $U_1$ على $5$ كريات منها: كريتان حمراوان وثلاث كريات خضراء، ويحتوي صندوق $U_2$ على $5$ كريات منها: ثلاث كريات حمراء وكريتان خضراوان. (جميع الكريات متماثلة ولا نفرق بينها باللمس)

نسحب عشوائيًا $3$ كريات بالكيفية التالية: نقوم بسحب كرية واحدة من $U_1$ ونسجل لونها.
\begin{itemize}
\item إذا كانت الكرية المسحوبة حمراء نعيدها إلى $U_1$ ونسحب منه عشوائيًا كريتين في آن واحد.
\item إذا كانت الكرية المسحوبة خضراء نضعها في $U_2$ ونسحب منه عشوائيًا كريتين في آن واحد.
\end{itemize}
نعتبر الحوادث: $R$: «الحصول على كرية حمراء»، \ $V$: «الحصول على كرية خضراء»،

$A$: «الحصول على $3$ كريات من نفس اللون»، \ $B$: «الحصول على كرية خضراء على الأقل».

\begin{center}\begin{LTR}\begin{tikzpicture}[bsT,level 1/.style={sibling distance=4.8cm},level 2/.style={sibling distance=1.6cm}]
\node[bsev]{}
 child{node[bsev]{$R$}
   child{node[bslf]{$RR$} edge from parent node[bsp]{$\dfrac{1}{10}$}}
   child{node[bslf]{$RV$} edge from parent node[bsp]{$\ldots$}}
   child{node[bslf]{$VV$} edge from parent node[bsp]{$\ldots$}}
   edge from parent node[bsp]{$\ldots$}}
 child{node[bsev]{$V$}
   child{node[bslf]{$RR$} edge from parent node[bsp]{$\dfrac{1}{5}$}}
   child{node[bslf]{$RV$} edge from parent node[bsp]{$\ldots$}}
   child{node[bslf]{$VV$} edge from parent node[bsp]{$\ldots$}}
   edge from parent node[bsp]{$\ldots$}};
\end{tikzpicture}\end{LTR}\end{center}
\QQ{1} أنقل وأكمل شجرة الاحتمالات المقابلة.
\QQ{2} أحسب احتمالي الحادثتين $A$ و$B$.
\QQ{3} $X$ المتغير العشوائي الذي يرفق بكل عملية سحب لثلاث كريات بالكيفية السابقة عدد الكريات الحمراء المسحوبة.

عيّن قانون احتمال المتغير العشوائي $X$ ثم أحسب أمله الرياضياتي.
\end{tamrin}

%======================== 4 ========================
\begin{tamrin}{بكالوريا 2025 -- الموضوع الثاني -- علوم تجريبية}
\textbf{I)} يحتوي صندوق $U_1$ على $6$ كريات متماثلة لا نفرق بينها باللمس، منها: $3$ كريات خضراء، وكريتان حمراوان، وكرية بيضاء واحدة.

نسحب عشوائيًا من الصندوق $U_1$ كريتين في آن واحد، ونعتبر الحادثتين:

$E$: «الحصول على كرية حمراء واحدة فقط»، \quad $F$: «الحصول على كريتين من نفس اللون».
\QQ{1} أحسب $P(E)$ احتمال الحادثة $E$.
\QQ{2} بيّن أن $P(F)$ احتمال الحادثة $F$ يساوي $\dfrac{4}{15}$، ثم استنتج احتمال الحصول على كريتين من لونين مختلفين.

\textbf{II)} يحتوي صندوق آخر $U_2$ على $6$ كريات متماثلة لا نفرق بينها باللمس، منها: $4$ كريات خضراء وكريتان حمراوان.
نرمز بـ $V$، $R$، $B$ إلى: خضراء، حمراء، بيضاء على الترتيب، ونسحب عشوائيًا من الصندوق $U_2$ كرية واحدة ونسجل لونها.
\begin{itemize}
\item إذا كانت الكرية المسحوبة خضراء، نسحب كرية أخرى من $U_2$ دون إرجاع الكرية الأولى.
\item وإذا كانت الكرية المسحوبة حمراء، نسحب كرية واحدة من $U_1$.
\end{itemize}
\begin{center}\begin{LTR}\begin{tikzpicture}[bsT,level 1/.style={sibling distance=4.4cm},level 2/.style={sibling distance=1.6cm}]
\node[bsev]{}
 child{node[bsev]{$V$}
   child{node[bsev]{$V$} edge from parent node[bsp]{$\ldots$}}
   child{node[bsev]{$R$} edge from parent node[bsp]{$\ldots$}}
   edge from parent node[bsp]{$\ldots$}}
 child{node[bsev]{$R$}
   child{node[bsev]{$V$} edge from parent node[bsp]{$\ldots$}}
   child{node[bsev]{$R$} edge from parent node[bsp]{$\ldots$}}
   child{node[bsev]{$B$} edge from parent node[bsp]{$\ldots$}}
   edge from parent node[bsp]{$\ldots$}};
\end{tikzpicture}\end{LTR}\end{center}
\QQ{1} أنقل واملأ شجرة الاحتمالات المقابلة.
\QQ{2}
\begin{sq}
\qa ما احتمال الحصول على كرية حمراء في السحب الثاني؟
\qb بيّن أن احتمال الحصول على كرية خضراء في السحب الأول علمًا أن الكرية المسحوبة الثانية حمراء يساوي $\dfrac{12}{17}$.
\end{sq}
\end{tamrin}

%======================== 5 ========================
\begin{tamrin}{بكالوريا 2024 -- الموضوع الأول -- علوم تجريبية}
يحتوي كيس على $11$ كرية متماثلة لا نفرق بينها باللمس موزعة كما يلي: كريتان بيضاوان مرقمتان بـ: $1$، $3$، وأربع كريات حمراء مرقمة بـ: $0$، $1$، $1$، $3$، وخمس كريات خضراء مرقمة بـ: $0$، $1$، $1$، $3$، $4$.

\textbf{I)} نسحب عشوائيًا وفي آن واحد $3$ كريات من الكيس، ونعتبر الأحداث الآتية:

$A$: «الحصول على $3$ كريات من نفس اللون»، \quad $B$: «الحصول على $3$ كريات جداء أرقامها عدد فردي»،

$C$: «الحصول على كريات جداء أرقامها عدد زوجي».
\QQ{1}
\begin{sq}
\qa أحسب $P(A)$ احتمال الحادثة $A$ وبيّن أن: $P(B)=\dfrac{56}{165}$ ثم استنتج $P(C)$.
\qb أحسب الاحتمال الشرطي $P_A(B)$.
\end{sq}
\QQ{2} $X$ المتغير العشوائي الذي يرفق بكل عملية سحب لثلاث كريات، عدد الكريات التي تحمل رقمًا زوجيًا.
\begin{sq}
\qa عيّن قانون الاحتمال للمتغير العشوائي $X$، ثم أحسب أمله الرياضياتي $E(X)$.
\qb أحسب احتمال الحادثة $(X>1)$.
\end{sq}
\textbf{II)} نسحب الآن من الكيس عشوائيًا $3$ كريات على التوالي وبدون إرجاع.

أحسب احتمال الحادثة $D$: «الحصول على $3$ كريات جداء أرقامها معدوم».
\end{tamrin}

%======================== 6 ========================
\begin{tamrin}{بكالوريا 2024 -- الموضوع الثاني -- علوم تجريبية}
يحتوي كيس على $5$ قطع كهربائية غير متمايزة ولا نفرق بينها باللمس، منها $3$ قطع سليمة وقطعتان غير سليمتين.
نرمز إلى القطعة السليمة بالرمز $S$ وإلى القطعة غير السليمة بالرمز $\overline{S}$.

نسحب عشوائيًا من الكيس $3$ قطع على التوالي مع الإرجاع، ونعتبر الحوادث:

$A$: «القطعة الأولى المسحوبة سليمة»، \quad $B$: «سحب قطعة واحدة فقط سليمة»،

$C$: «القطعة الثالثة المسحوبة سليمة».
\QQ{1} شكّل شجرة الاحتمالات التي تنمذج هذه التجربة.
\QQ{2} أحسب $P(A)$ و$P(B)$ احتمالي الحادثتين $A$ و$B$ ثم بيّن أن: $P(C)=\dfrac35$.
\QQ{3} أحسب الاحتمال الشرطي $P_C(A)$. هل الحادثتان $A$ و$C$ مستقلتان؟
\QQ{4} نرفق بكل قطعة سليمة العدد $10$ وبكل قطعة غير سليمة العدد $-10$، ونعتبر $X$ المتغير العشوائي الذي يرفق بكل عملية سحب لثلاث قطع مجموع الأعداد المرفقة بها.
\begin{sq}
\qa برّر أن قيم المتغير العشوائي $X$ هي: $-30$، $-10$، $10$، $30$.
\qb عيّن قانون الاحتمال للمتغير العشوائي $X$ ثم أحسب أمله الرياضياتي $E(X)$.
\end{sq}
\end{tamrin}

%======================== 7 ========================
\begin{tamrin}{بكالوريا 2023 -- الموضوع الأول -- علوم تجريبية}
يحتوي صندوق $U_1$ على $5$ كريات تحمل الأرقام $1$، $1$، $1$، $2$، $3$، ويحتوي صندوق $U_2$ على $4$ كريات تحمل الأرقام $1$، $1$، $2$، $2$. (كل الكريات متماثلة ولا نفرق بينها عند اللمس)

نختار عشوائيًا أحد الصندوقين ونسحب منه عشوائيًا كريتين في آن واحد.
\QQ{1} نعتبر الأحداث: $A$: «سحب كريتين تحملان رقمين فرديين»، \ $B$: «سحب كريتين تحملان رقمين زوجيين»،

$C$: «سحب كريتين إحداهما تحمل رقمًا فرديًا والأخرى تحمل رقمًا زوجيًا».
\begin{sq}
\qa أنجز الشجرة التي تنمذج هذه التجربة.
\qb بيّن أن: $P(A)=\dfrac{23}{60}$ و$P(B)=\dfrac{1}{12}$ ثم أحسب $P(C)$.
\end{sq}
\QQ{2} نفرغ محتوى الصندوقين $U_1$ و$U_2$ في صندوق جديد $U_3$ ثم نسحب منه عشوائيًا كريتين في آن واحد.

$X$ المتغير العشوائي الذي يرفق بكل عملية سحب لكريتين جداء الرقمين المسجلين عليهما.
\begin{sq}
\qa برّر أن مجموعة قيم المتغير العشوائي $X$ هي: $\{1;2;3;4;6\}$.
\qb عيّن قانون الاحتمال للمتغير العشوائي $X$ ثم أحسب أمله الرياضياتي $E(X)$.
\end{sq}
\end{tamrin}

%======================== 8 ========================
\begin{tamrin}{بكالوريا 2023 -- الموضوع الثاني -- علوم تجريبية}
يحتوي كيس على $10$ كريات متماثلة ولا نفرق بينها باللمس، موزعة كما يلي: ثلاث كريات بيضاء مرقمة بـ: $1$، $1$، $2$، وثلاث كريات حمراء مرقمة بـ: $1$، $2$، $2$، وأربع كريات خضراء مرقمة بـ: $1$، $2$، $2$، $2$.

نسحب عشوائيًا وفي آن واحد كريتين من الكيس ونعتبر الأحداث الآتية:

$A$: «الحصول على كريتين من نفس اللون»، \quad $B$: «الحصول على كرية خضراء على الأقل»،

$C$: «الحصول على كريتين تحملان رقمين زوجيين».
\QQ{1}
\begin{sq}
\qa بيّن أن احتمال الحدث $A$ يساوي $\dfrac{4}{15}$ وأن احتمال الحدث $B$ يساوي $\dfrac23$.
\qb أحسب الاحتمالين $P(C)$ و$P(A\cap C)$، هل الحدثان $A$ و$C$ مستقلان؟
\qc استنتج احتمال الحصول على كريتين من نفس اللون علمًا أنهما تحملان رقمين زوجيين.
\end{sq}
\QQ{2} نعتبر المتغير العشوائي $X$ الذي يرفق بكل عملية سحب لكريتين مجموع الرقمين المسجلين عليهما.
\begin{sq}
\qa برّر أن مجموعة قيم المتغير العشوائي $X$ هي: $\{2;3;4\}$.
\qb عيّن قانون احتمال المتغير العشوائي $X$ ثم أحسب أمله الرياضياتي $E(X)$.
\end{sq}
\end{tamrin}

%======================== 9 ========================
\begin{tamrin}{بكالوريا 2021 -- الموضوع الأول -- علوم تجريبية}
يُراد تشكيل بطريقة عشوائية لجنة تتكون من عضوين من بين ثلاثة رجال $H_1$، $H_2$ و$H_3$ وامرأتين $F_1$ و$F_2$.

نعتبر الأحداث $A$، $B$ و$C$ حيث: \ $A$: «عضوا اللجنة من نفس الجنس»،

$B$: «عضوا اللجنة من جنسين مختلفين»، \quad $C$: «$H_1$ عضو في اللجنة».
\QQ{1}
\begin{sq}
\qa أحسب $P(A)$ و$P(B)$ احتمال $A$ و$B$ على الترتيب.
\qb بيّن أن $P(C)$ احتمال الحدث $C$ يساوي $\dfrac25$.
\end{sq}
\QQ{2} المتغير العشوائي $X$ يرفق بكل إمكانية اختيار لعضوين عدد الرجال في اللجنة.
\begin{sq}
\qa برّر أن مجموعة قيم $X$ هي: $\{0;1;2\}$.
\qb عيّن قانون احتمال المتغير العشوائي $X$ وأحسب أمله الرياضياتي $E(X)$.
\end{sq}
\end{tamrin}

%======================== 10 ========================
\begin{tamrin}{بكالوريا 2021 -- الموضوع الثاني -- علوم تجريبية}
صندوق به $9$ بطاقات متماثلة لا نفرق بينها باللمس، مكتوب على كل منها سؤال واحد، منها ثلاثة أسئلة في الهندسة مرقمة بـ: $1$، $2$ و$3$، أربعة أسئلة في الجبر مرقمة بـ: $1$، $2$، $3$ و$4$، وسؤالان في التحليل مرقمان بـ: $1$ و$2$.

نسحب عشوائيًا بطاقة واحدة من الصندوق ونعتبر الأحداث التالية:

$A$: «سحب سؤال في الهندسة»، \ $B$: «سحب سؤال في التحليل»، \ $C$: «سحب سؤال في الجبر يحمل رقمًا زوجيًا».
\QQ{1} أحسب $P(A)$، $P(B)$ و$P(C)$ احتمالات الأحداث $A$، $B$ و$C$ على الترتيب.
\QQ{2} أحسب احتمال سحب سؤال رقمه مختلف عن $1$.
\QQ{3} المتغير العشوائي $X$ يرفق بكل بطاقة مسحوبة رقم السؤال المسجل عليها.
\begin{sq}
\qa برّر أن مجموعة قيم $X$ هي: $\{1;2;3;4\}$.
\qb عيّن قانون احتمال المتغير العشوائي $X$ ثم أحسب $E(X)$ أمله الرياضياتي.
\qc استنتج قيمة $E(2021X+1442)$.
\end{sq}
\end{tamrin}

%======================== 11 ========================
\begin{tamrin}{بكالوريا 2020 -- الموضوع الأول -- علوم تجريبية}
يحتوي وعاء $U$ على $4$ كريات حمراء و$6$ سوداء، ويحتوي وعاء $V$ على $5$ كريات حمراء و$3$ سوداء، وكل الكريات متماثلة ولا نفرق بينها عند اللمس.

نسحب عشوائيًا كريتين في آن واحد من أحد الوعاءين بالكيفية التالية: نقوم بسحب بطاقة واحدة عشوائيًا من كيس يحتوي على $6$ بطاقات متماثلة ومرقمة من $1$ إلى $6$؛ إذا تحصلنا على أحد الرقمين $3$ أو $5$ نسحب الكريتين من $U$، وفي باقي الحالات نسحب الكريتين من $V$.

نسمي $A$ الحدث: «الحصول على أحد الرقمين $3$ أو $5$»، ونسمي $M$ الحدث: «الحصول على كريتين من نفس اللون».
\begin{center}\begin{LTR}\begin{tikzpicture}[bsT,level 1/.style={sibling distance=2.8cm},level 2/.style={sibling distance=1.5cm}]
\node[bsev]{}
 child{node[bsev]{$A$}
   child{node[bsev]{$M$} edge from parent node[bsp]{$\ldots$}}
   child{node[bsev]{$\overline{M}$} edge from parent node[bsp]{$\ldots$}}
   edge from parent node[bsp]{$\ldots$}}
 child{node[bsev]{$\overline{A}$}
   child{node[bsev]{$M$} edge from parent node[bsp]{$\dfrac{13}{28}$}}
   child{node[bsev]{$\overline{M}$} edge from parent node[bsp]{$\ldots$}}
   edge from parent node[bsp]{$\ldots$}};
\end{tikzpicture}\end{LTR}\end{center}
\QQ{1} تحقق أن احتمال السحب من الوعاء $V$ هو $P(\overline{A})=\dfrac23$.
\QQ{2} علمًا أن الكريتين المسحوبتين من $U$، بيّن أن احتمال أن تكونا من نفس اللون هو $\dfrac{7}{15}$.
\QQ{3} أنقل شجرة الاحتمالات المقابلة ثم أكملها واستنتج $P(M)$.
\QQ{4} أحسب $P_{\overline{M}}(A)$ احتمال السحب من الوعاء $U$ علمًا أن الكريتين المسحوبتين مختلفتا اللون.
\end{tamrin}

%======================== 12 ========================
\begin{tamrin}{بكالوريا 2020 -- الموضوع الثاني -- علوم تجريبية}
كيس به ثلاث كريات بيضاء وكريتان حمراوان لا نميز بينها عند اللمس، نسحب عشوائيًا كريتين على التوالي من الكيس بالكيفية التالية: إذا كانت الكرية المسحوبة بيضاء نعيدها إلى الكيس، وإذا كانت حمراء لا نعيدها إلى الكيس.
\begin{center}\begin{LTR}\begin{tikzpicture}[bsT,level 1/.style={sibling distance=2.8cm},level 2/.style={sibling distance=1.5cm}]
\node[bsev]{}
 child{node[bsev]{$B$}
   child{node[bsev]{$B$} edge from parent node[bsp]{$\ldots$}}
   child{node[bsev]{$R$} edge from parent node[bsp]{$\dfrac25$}}
   edge from parent node[bsp]{$\ldots$}}
 child{node[bsev]{$R$}
   child{node[bsev]{$B$} edge from parent node[bsp]{$\dfrac34$}}
   child{node[bsev]{$R$} edge from parent node[bsp]{$\ldots$}}
   edge from parent node[bsp]{$\ldots$}};
\end{tikzpicture}\end{LTR}\end{center}
\QQ{1}
\begin{sq}
\qa أنقل شجرة الاحتمالات المقابلة ثم أكملها. ($B$ يرمز إلى الحصول على كرية بيضاء و$R$ إلى الحصول على كرية حمراء.)
\qb أحسب احتمال أن تكون الكرية المسحوبة الثانية حمراء.
\end{sq}
\QQ{2} ليكن $X$ المتغير العشوائي الذي يرفق بكل سحب لكريتين عدد الكريات الحمراء المسحوبة.
\begin{sq}
\qa عيّن مجموعة قيم المتغير العشوائي $X$.
\qb بيّن أن: $P(X=1)=\dfrac{27}{50}$، ثم عرّف قانون احتمال المتغير العشوائي $X$.
\qc أحسب $E(X)$ الأمل الرياضياتي للمتغير العشوائي $X$.
\end{sq}
\end{tamrin}

%======================== 13 ========================
\begin{tamrin}{بكالوريا 2019 -- الموضوع الأول -- علوم تجريبية}
يحتوي كيس على خمس كريات حمراء، منها أربع كريات تحمل الرقم $1$ وكرية واحدة تحمل الرقم $2$، وسبع كريات خضراء، منها أربع كريات تحمل الرقم $1$ وثلاث كريات تحمل الرقم $2$. (كل الكريات متماثلة لا نفرق بينها عند اللمس)

نسحب عشوائيًا كريتين من الكيس في آن واحد ونعتبر الحادثتين $A$ و$B$ حيث:

$A$: «سحب كريتين من نفس اللون»، \quad $B$: «سحب كريتين تحملان نفس الرقم».
\QQ{1} بيّن أن احتمال الحادثة $A$ هو $P(A)=\dfrac{31}{66}$ وأحسب احتمال الحادثة $B$.
\QQ{2} علمًا أن الكريتين المسحوبتين من نفس اللون، ما احتمال أن تحملا نفس الرقم؟
\QQ{3} ليكن $X$ المتغير العشوائي الذي يرفق بكل عملية سحب عدد الكريات الحمراء المتبقية في الكيس.

عرّف قانون الاحتمال للمتغير العشوائي وأحسب أمله الرياضياتي.
\end{tamrin}

%======================== 14 ========================
\begin{tamrin}{بكالوريا 2019 -- الموضوع الثاني -- علوم تجريبية}
يحتوي صندوق على $10$ كريات لا نفرق بينها عند اللمس، منها كريتان تحملان الرقم $0$، وثلاث تحمل الرقم $1$، والكريات الأخرى تحمل الرقم $2$. نسحب عشوائيًا وفي آن واحد ثلاث كريات من الصندوق.

ليكن $X$ المتغير العشوائي الذي يرفق بكل سحب جداء الأرقام المسجلة على الكريات المسحوبة.
\QQ{1} عرّف قانون الاحتمال للمتغير العشوائي $X$، ثم أحسب أمله الرياضياتي $E(X)$.
\QQ{2} بيّن أن احتمال الحصول على ثلاث كريات كل منها تحمل رقمًا زوجيًا هو $\dfrac{7}{24}$.
\QQ{3} نسحب الآن من الصندوق كريتين على التوالي دون إرجاع.

ما احتمال الحصول على كريتين تحملان رقمين مجموعهما فردي علمًا أن جداءهما زوجي؟
\end{tamrin}

%======================== 15 ========================
\begin{tamrin}{بكالوريا 2018 -- الموضوع الأول -- علوم تجريبية}
يحتوي صندوق $10$ كريات متماثلة لا نفرق بينها باللمس، منها $4$ كريات بيضاء مرقمة بـ: $1$، $2$، $2$، $3$، وثلاث كريات حمراء مرقمة بـ: $2$، $2$، $3$، وثلاث كريات خضراء مرقمة بـ: $2$، $3$، $3$.
نسحب عشوائيًا وفي آن واحد $3$ كريات من هذا الصندوق.

نعتبر الحادثتين: $A$: «الكريات الثلاث المسحوبة تحمل ألوان العلم الوطني»، و$B$: «الكريات الثلاث المسحوبة لها نفس الرقم».
\QQ{1}
\begin{sq}
\qa أحسب $P(A)$ و$P(B)$ احتمالي الحادثتين $A$ و$B$ على الترتيب.
\qb بيّن أن: $P(A\cap B)=\dfrac{1}{20}$ ثم استنتج $P_A(B)$ و$P(A\cup B)$.
\end{sq}
\QQ{2} ليكن $X$ المتغير العشوائي الذي يرفق بكل نتيجة عملية سحب عدد الكريات التي تحمل رقمًا فرديًا.

عرّف قانون الاحتمال للمتغير العشوائي $X$ وأحسب أمله الرياضياتي $E(X)$.
\end{tamrin}

\section{تمارين البكالوريا -- شعبة تقني رياضي}
% ===================== تمارين البكالوريا -- تقني رياضي =====================
\tikzset{btree/.style={grow=right,level distance=3.2cm,
  ev/.style={circle,draw=cDef,fill=cDef!12,thick,minimum size=7mm,inner sep=1pt,font=\small},
  edge from parent/.style={draw,thick,cDef!70!black},
  pu/.style={midway,above,sloped,inner sep=1pt,font=\small},
  pd/.style={midway,below,sloped,inner sep=1pt,font=\small},
  lf/.style={label={[label distance=5mm,font=\small]right:{##1}}}}}

% ===================== التمرين 1 =====================
\begin{tamrin}{بكالوريا 2026 -- الموضوع الأول -- تقني رياضي}
يحتوي كيس على $10$ كريات متماثلة ولا نفرق بينها باللمس، منها: كريتان بيضاوان مرقمتان بـ: $0$، $2$ وأربع كريات حمراء مرقمة بـ: $0$، $1$، $2$، $3$ وأربع كريات خضراء مرقمة بـ: $1$، $2$، $3$، $4$.

نسحب عشوائيا من الكيس ثلاث كريات في آن واحد ونعتبر الحوادث التالية:
\begin{itemize}
\item $A$: «سحب كرية حمراء على الأقل»، \quad $B$: «سحب ثلاث كريات جداء أرقامها معدوم»،
\item $C$: «سحب ثلاث كريات من نفس اللون»، \quad $D$: «سحب ثلاث كريات مجموع أرقامها أكبر من أو يساوي $8$».
\end{itemize}
\QQ{1}
\begin{sq}
\qa أحسب $P(A)$، $P(B)$ و $P(C)$ احتمالات الحوادث $A$، $B$ و $C$ على الترتيب.
\qb بيّن أن: $P(D)=\dfrac{17}{120}$ وأن: $P(C\cap D)=\dfrac{1}{60}$.
\qc استنتج حساب احتمال الحادثة $E$: «سحب ثلاث كريات مجموع أرقامها أصغر من أو يساوي $7$».
\end{sq}
\QQ{2} $X$ المتغيّر العشوائي الذي يرفق بكل عملية سحب لثلاث كريات عدد الكريات التي تحمل رقما فرديا.

عيّن قانون الاحتمال للمتغيّر العشوائي $X$، ثم أحسب أمله الرياضياتي $E(X)$.
\end{tamrin}

% ===================== التمرين 2 =====================
\begin{tamrin}{بكالوريا 2026 -- الموضوع الثاني -- تقني رياضي}
يحتوي كيس على $11$ كرية متماثلة ولا نفرق بينها باللمس، منها: كريتان بيضاوان، خمس كريات حمراء، ثلاث كريات خضراء وكرية واحدة سوداء.
\QQ{1} نسحب عشوائيا من الكيس ثلاث كريات على التوالي دون إرجاع ونعتبر الحوادث التالية:
\begin{itemize}
\item $A$: «سحب ثلاث كريات من نفس اللون»، \quad $B$: «سحب كرية واحدة بيضاء فقط»،
\item $C$: «سحب كرية حمراء على الأكثر»، \quad $D$: «سحب كريتين حمراوين على الأقل».
\end{itemize}
\begin{sq}
\qa أحسب $P(A)$، $P(B)$، $P(C)$ و $P(D)$ احتمالات الحوادث $A$، $B$، $C$ و $D$ على الترتيب.
\qb بيّن أن: $P(A\cap D)=\dfrac{2}{33}$ ثم استنتج حساب $P(A\cup D)$.
\end{sq}
\QQ{2} نفرض أن عدد الكريات السوداء في الكيس هو $n$ حيث $n\ge 2$ ونسحب منه كريتين على التوالي مع الإرجاع.

عيّن قيمة $n$ حتى يكون احتمال سحب كريتين سوداوين يساوي $\dfrac{1}{4}$.
\end{tamrin}

% ===================== التمرين 3 =====================
\begin{tamrin}{بكالوريا 2025 -- الموضوع الأول -- تقني رياضي}
يحتوي صندوق على $10$ كريات متماثلة ولا نفرق بينها باللمس، منها: كريتان خضراوان مرقمتان بـ: $1$، $2$ وثلاث كريات بيضاء مرقمة بـ: $1$، $1$، $2$ وخمس كريات حمراء مرقمة بـ: $1$، $1$، $2$، $2$، $3$.

\textbf{I)} نسحب عشوائيا من الصندوق $3$ كريات في آن واحد ونعتبر الحوادث التالية:
\begin{itemize}
\item $A$: «الحصول على ثلاث كريات مختلفة الألوان مثنى مثنى»، \quad $B$: «الحصول على ثلاث كريات تحمل نفس الرقم»،
\item $C$: «الحصول على ثلاث كريات مجموع أرقامها $6$».
\end{itemize}
\QQ{1}
\begin{sq}
\qa أحسب $P(A)$، $P(B)$ و $P(C)$ احتمالات الحوادث $A$، $B$ و $C$ على الترتيب.
\qb أحسب $P(A\cap B)$ ثم استنتج حساب $P(A\cup B)$.
\end{sq}
\QQ{2} $X$ المتغيّر العشوائي الذي يرفق بكل عملية سحب لـ $3$ كريات بالكيفية السابقة عدد ألوان الكريات المسحوبة.

عيّن قانون الاحتمال للمتغيّر العشوائي $X$، ثم أحسب أمله الرياضياتي $E(X)$.

\textbf{II)} نسحب عشوائيا من الصندوق $3$ كريات على التوالي دون إرجاع. ما احتمال ألا يبقى في الصندوق أية كرية خضراء؟
\end{tamrin}

% ===================== التمرين 4 =====================
\begin{tamrin}{بكالوريا 2025 -- الموضوع الثاني -- تقني رياضي}
يحتوي صندوق على $12$ كرية منها: $3$ كريات خضراء و $4$ كريات بيضاء و $5$ كريات حمراء. كل الكريات متماثلة ولا نفرق بينها باللمس.

\textbf{I)} نسحب عشوائيا من الصندوق $4$ كريات في آن واحد ونعتبر الحوادث التالية:
\begin{itemize}
\item $A$: «الحصول على $3$ كريات بالضبط من نفس اللون»، \quad $B$: «لا نحصل على أية كرية بيضاء»،
\item $C$: «الحصول على كرية واحدة بيضاء على الأقل».
\end{itemize}
\QQ{1}
\begin{sq}
\qa أحسب $P(A)$، $P(B)$ و $P(C)$ احتمالات الحوادث $A$، $B$ و $C$ على الترتيب.
\qb أحسب $P(A\cap B)$ ثم استنتج حساب $P(A\cup B)$.
\end{sq}
\QQ{2} $X$ المتغيّر العشوائي الذي يرفق بكل عملية سحب لـ $4$ كريات بالكيفية السابقة عدد الكريات الخضراء المسحوبة.

عيّن قانون الاحتمال للمتغيّر العشوائي $X$، ثم أحسب أمله الرياضياتي $E(X)$.

\textbf{II)} نسحب عشوائيا من الصندوق $4$ كريات على التوالي دون إرجاع. ما احتمال ألا يبقى في الصندوق أية كرية خضراء؟
\end{tamrin}

% ===================== التمرين 5 =====================
\begin{tamrin}{بكالوريا 2024 -- الموضوع الأول -- تقني رياضي}
يحتوي صندوق على $9$ كريات متماثلة ولا نفرق بينها باللمس، منها: $4$ بيضاء و $3$ حمراء و $2$ خضراوان.

\textbf{I)} نسحب عشوائيا من الصندوق $3$ كريات في آن واحد.
\QQ{1} أحسب احتمال كل من الحوادث الآتية:
\begin{itemize}
\item $A$: «الحصول على ثلاث كريات من نفس اللون»، \quad $B$: «الحصول على الألوان الثلاثة»،
\item $C$: «الحصول على كرية بيضاء على الأقل».
\end{itemize}
\QQ{2} $X$ المتغيّر العشوائي الذي يرفق بكل عملية سحب لثلاث كريات عدد الألوان المتحصل عليها.
\begin{sq}
\qa عيّن قانون الاحتمال للمتغيّر العشوائي $X$، ثم أحسب أمله الرياضياتي $E(X)$.
\qb أحسب $E(84X+1837)$.
\end{sq}
\textbf{II)} نضيف الآن إلى الصندوق كرية واحدة سوداء ثم نسحب منه عشوائيا $4$ كريات على التوالي دون إرجاع.

بيّن أن احتمال الحادثة $D$: «الحصول على الألوان الأربعة» هو $\dfrac{4}{35}$.
\end{tamrin}

% ===================== التمرين 6 =====================
\begin{tamrin}{بكالوريا 2024 -- الموضوع الثاني -- تقني رياضي}
يحتوي كيس على $10$ كريات متماثلة ولا نفرق بينها باللمس منها: كريتان تحملان الرقم $0$، ثلاث كريات تحمل الرقم $2$، كرية تحمل الرقم $3$ وأربع كريات تحمل الرقم $4$. نسحب عشوائيا من الكيس ثلاث كريات في آن واحد.
\QQ{1} أحسب احتمال كل من الحوادث الآتية:
\begin{itemize}
\item $A$: «مجموع الأعداد التي تحملها الكريات المسحوبة يساوي $12$»،
\item $B$: «الحصول على ثلاثة أعداد أولية»،
\item $C$: «جداء الأعداد التي تحملها الكريات المسحوبة معدوم».
\end{itemize}
\QQ{2} $X$ المتغيّر العشوائي الذي يرفق بكل عملية سحب لثلاث كريات عدد الأعداد الأولية المتحصل عليها.
\begin{sq}
\qa عيّن قانون الاحتمال للمتغيّر العشوائي $X$ ثم أحسب أمله الرياضياتي $E(X)$.
\qb أحسب احتمال الحادثة $(X^2>e)$.
\end{sq}
\end{tamrin}

% ===================== التمرين 7 =====================
\begin{tamrin}{بكالوريا 2023 -- الموضوع الأول -- تقني رياضي}
يحتوي كيس على $8$ كريات متماثلة ولا نفرق بينها باللمس، موزعة كما يلي: ثلاث كريات بيضاء مرقمة بـ: $0$، $1$، $1$ وثلاث كريات حمراء مرقمة بـ: $1$، $1$، $2$ وكريتين خضراوين مرقمتين بـ: $1$، $2$.

نسحب عشوائيا وفي آن واحد كريتين من الكيس ونعتبر الأحداث $A$، $B$ و $C$ الآتية:
\begin{itemize}
\item $A$: «الحصول على كريتين من نفس اللون»، \quad $B$: «الحصول على كرية حمراء على الأقل»،
\item $C$: «الحصول على كريتين تحملان رقمين مجموعهما يساوي $3$».
\end{itemize}
\QQ{1}
\begin{sq}
\qa بيّن أن احتمال الحدث $A$ يساوي $\dfrac{1}{4}$، وأن احتمال الحدث $B$ يساوي $\dfrac{9}{14}$.
\qb أحسب الاحتمال $P(C)$.
\end{sq}
\QQ{2} نعتبر المتغيّر العشوائي $X$ الذي يرفق بكل عملية سحب لكريتين مجموع الرقمين المسجلين عليهما.
\begin{sq}
\qa برّر أن مجموعة قيم المتغيّر العشوائي $X$ هي: $\{1;2;3;4\}$.
\qb عيّن قانون احتمال المتغيّر العشوائي $X$ ثم أحسب أمله الرياضياتي $E(X)$.
\end{sq}
\end{tamrin}

% ===================== التمرين 8 =====================
\begin{tamrin}{بكالوريا 2023 -- الموضوع الثاني -- تقني رياضي}
يحتوي كيس على $11$ كرية متماثلة ولا نفرق بينها باللمس، موزعة كما يلي: ثلاث كريات تحمل الرقم $0$، ثلاث كريات تحمل الرقم $1$ وخمس كريات تحمل الرقم $2$.

نسحب عشوائيا وفي آن واحد كريتين من الكيس ونعتبر الأحداث $A$، $B$ و $C$ الآتية:
\begin{itemize}
\item $A$: «الحصول على كريتين رقم كل منهما عدد أولي»، \quad $B$: «الحصول على كرية واحدة تحمل رقما فرديا»،
\item $C$: «الحصول على كريتين جداء رقميهما معدوم».
\end{itemize}
\QQ{1}
\begin{sq}
\qa بيّن أن احتمال الحدث $A$ يساوي $\dfrac{2}{11}$ وأن احتمال الحدث $B$ يساوي $\dfrac{24}{55}$.
\qb أحسب الاحتمال $P(C)$.
\end{sq}
\QQ{2} نعتبر المتغيّر العشوائي $X$ الذي يرفق بكل عملية سحب لكريتين جداء الرقمين المسجلين عليهما.
\begin{sq}
\qa برّر أن مجموعة قيم المتغيّر العشوائي $X$ هي: $\{0;1;2;4\}$.
\qb عيّن قانون احتمال المتغيّر العشوائي $X$ ثم أحسب أمله الرياضياتي $E(X)$.
\qc أحسب احتمال الحدث: «$e^{X+6}<2023$».
\end{sq}
\end{tamrin}

% ===================== التمرين 9 =====================
\begin{tamrin}{بكالوريا 2020 -- الموضوع الأول -- تقني رياضي}
يحتوي كيس على أربع كريات حمراء مرقمة بـ: $2$، $2$، $2$، $2$ وثلاث كريات خضراء مرقمة بـ: $2$، $3$، $3$. الكريات لا نفرق بينها باللمس، نسحب عشوائيا وفي آن واحد كريتين من هذا الكيس.
\QQ{1} نعتبر الحدثين: $A$: «الحصول على كريتين تحملان نفس الرقم» و $B$: «الحصول على كريتين مختلفتين في اللون».
\begin{sq}
\qa أحسب احتمال كل من الحدثين $A$ و $B$.
\qb بيّن أن احتمال الحصول على كريتين تحملان نفس الرقم ومختلفتين في اللون يساوي $\dfrac{4}{21}$.
\qc استنتج احتمال الحصول على كريتين تحملان نفس الرقم أو مختلفتين في اللون.
\end{sq}
\QQ{2} ليكن $X$ المتغيّر العشوائي الذي يرفق بكل سحب جداء الرقمين الظاهرين على الكريتين المسحوبتين. عرّف قانون احتمال المتغيّر العشوائي $X$.
\QQ{3} في لعبة، يقوم لاعب بسحب كريتين: إذا كان جداء رقميهما $4$ يربح $x^2$ دينار، إذا كان جداء رقميهما $6$ يخسر $y^2$ دينار، وإذا كان جداء رقميهما $9$ يخسر $130$ دينار ($x$ و $y$ عددان طبيعيان غير معدومين).

عيّن قيمة كل من $x$ و $y$ حتى تكون هذه اللعبة عادلة.
\end{tamrin}

% ===================== التمرين 10 =====================
\begin{tamrin}{بكالوريا 2020 -- الموضوع الثاني -- تقني رياضي}
يحتوي كيس على كريتين خضراوين تحملان الرقمين $1$، $2$ وثلاث كريات حمراء تحمل الأرقام $1$، $2$، $2$ وأربع كريات بيضاء تحمل الأرقام $2$، $3$، $3$، $4$ (الكريات متماثلة لا نفرق بينها باللمس).

\textbf{I)} نسحب من هذا الكيس ثلاث كريات في آن واحد.
\QQ{1} أحسب احتمال كل من الحدثين $A$ و $B$ التاليين: $A$: «الحصول على ثلاث كريات من نفس اللون» و $B$: «الحصول على كرية بيضاء على الأقل».
\QQ{2} ليكن $X$ المتغيّر العشوائي الذي يرفق بكل سحب أكبر الأرقام المحصل عليها.
\begin{sq}
\qa بيّن أن $P(X=3)=\dfrac{3}{7}$، ثم عرّف قانون الاحتمال للمتغيّر العشوائي $X$.
\qb أحسب الأمل الرياضياتي للمتغيّر العشوائي $X$.
\end{sq}
\textbf{II)} نسحب الآن ثلاث كريات على التوالي دون إرجاع. ليكن $C$ الحدث «الحصول على ثلاثة أرقام جداؤها زوجي». أحسب احتمال $C$.
\end{tamrin}

% ===================== التمرين 11 =====================
\begin{tamrin}{بكالوريا 2019 -- الموضوع الأول -- تقني رياضي}
توجد إجابة صحيحة واحدة من بين الأجوبة المقترحة في كل حالة من الحالات التالية. اختر الإجابة الصحيحة مبررا اختيارك.

يحتوي كيس على ثلاث كريات بيضاء تحمل الأرقام $1$، $2$، $3$ وكريتين سوداوين تحملان الرقمين $1$، $2$ (الكريات لا نفرق بينها عند اللمس). نسحب من الكيس $3$ كريات عشوائيا وفي آن واحد.

$X$ المتغيّر العشوائي الذي يرفق بكل سحب عدد الكريات السوداء المسحوبة.
\QQ{1} قيم المتغيّر العشوائي $X$ هي:
\ABC{$\{1;2;3\}$}{$\{0;2;3\}$}{$\{0;1;2\}$}
\QQ{2} الأمل الرياضياتي $E(X)$ لـ $X$ هو:
\ABC{$E(X)=\dfrac45$}{$E(X)=\dfrac65$}{$E(X)=\dfrac{11}{10}$}
\QQ{3} احتمال «الحصول على كرية واحدة سوداء تحمل الرقم $1$ من الكريات المسحوبة» يساوي:
\ABC{$\dfrac{7}{10}$}{$\dfrac{9}{10}$}{$\dfrac35$}
\QQ{4} احتمال «باقي قسمة مجموع مربعات الأرقام التي تحملها الكريات المسحوبة على $13$ هو $1$» يساوي:
\ABC{$\dfrac25$}{$\dfrac{3}{10}$}{$\dfrac15$}
\end{tamrin}

% ===================== التمرين 12 =====================
\begin{tamrin}{بكالوريا 2019 -- الموضوع الثاني -- تقني رياضي}
يحتوي كيس على أربع كريات بيضاء تحمل الأرقام $1$، $2$، $3$، $4$ وثلاث كريات حمراء تحمل الأرقام $1$، $2$، $3$ وكريتين سوداوين تحملان الرقمين $1$، $2$ (كل الكريات متشابهة لا نفرق بينها عند اللمس). نسحب عشوائيا وفي آن واحد $3$ كريات من هذا الكيس.
\QQ{1} أحسب احتمال الأحداث التالية:
\begin{itemize}
\item $A$: «الحصول على كرية بيضاء واحدة».
\item $B$: «الحصول على كريتين بيضاوين على الأكثر».
\item $C$: «الحصول على ثلاث كريات تحمل أرقاما غير أولية».
\end{itemize}
\QQ{2} نعتبر المتغيّر العشوائي $X$ الذي يرفق بكل عملية سحب عدد الكريات التي تحمل أرقاما أولية.
\begin{sq}
\qa عيّن قيم المتغيّر العشوائي $X$، ثم عرّف قانون احتماله.
\qb أحسب $P(X^2-X\le0)$.
\end{sq}
\end{tamrin}

% ===================== التمرين 13 =====================
\begin{tamrin}{بكالوريا 2018 -- الموضوع الثاني -- تقني رياضي}
كيس به $7$ كريات متماثلة لا نفرق بينها باللمس، منها $3$ بيضاء و $4$ خضراء. نسحب عشوائيا وفي آن واحد كريتين من الكيس.

\textbf{I)}
\QQ{1} أحسب احتمال الحادثة $A$: «سحب كريتين مختلفتين في اللون».
\QQ{2} أحسب احتمال الحادثة $B$: «سحب كريتين من نفس اللون».

\textbf{II)} نقترح اللعبة التالية: للمشاركة يدفع اللاعب $\alpha$ (DA)، (حيث $\alpha$ عدد طبيعي معطى و DA تعني دينار جزائري).
فإذا سحب كريتين بيضاوين يتحصل على $100\,\text{DA}$، وإذا سحب كريتين مختلفتين في اللون يتحصل على $50\,\text{DA}$، وإذا سحب كريتين خضراوين يخسر ما دفعه. وليكن $X$ المتغيّر العشوائي الذي يمثل ربح أو خسارة اللاعب بدلالة $\alpha$.
\QQ{1} برّر أن قيم المتغيّر العشوائي هي $\{100-\alpha;\,50-\alpha;\,-\alpha\}$ ثم عرّف قانون احتماله.
\QQ{2} بيّن أن الأمل الرياضياتي للمتغيّر العشوائي $X$ بدلالة $\alpha$ هو: $E(X)=-\alpha+\dfrac{300}{7}$، ثم جد أكبر قيمة ممكنة لـ $\alpha$ حتى تكون اللعبة في صالح اللاعب.
\end{tamrin}

\section{تمارين البكالوريا -- شعبة رياضيات}
% ===== تمارين البكالوريا في الاحتمالات -- شعبة رياضيات =====
\tikzset{
  bmT/.style={grow'=right,level distance=3.4cm,edge from parent/.style={draw,thick,cDef!70!black}},
  bmS/.style={grow'=right,level distance=2.4cm,edge from parent/.style={draw,thick,gray!70!black}},
  ev/.style={circle,draw=cDef,fill=cDef!12,thick,minimum size=7mm,inner sep=1pt,font=\small},
  p/.style={midway,fill=white,inner sep=1.5pt,font=\footnotesize},pb/.style={pos=0.62,fill=white,inner sep=1pt,font=\footnotesize},
  pl/.style={font=\small,text=cSec,label distance=3mm},
  hd/.style={font=\small\bfseries,text=cSec}}

%======================================================================
\begin{tamrin}{بكالوريا 2026 -- الموضوع الأول -- شعبة رياضيات}
يحتوي كيس $U$ على $4$ كريات حمراء و$3$ كريات خضراء، ويحتوي كيس $V$ على كريتين حمراوين وكريتين خضراوين.
كل الكريات متماثلة ولا نفرق بينها باللمس.

نسحب عشوائيا كريتين في آن واحد من $U$؛ إذا كانتا حمراوين نضعهما في $V$ ثم نسحب منه عشوائيا كريتين في آن واحد، وإذا لم تكونا حمراوين معا نعيدهما إلى $U$ ثم نسحب من $V$ كريتين في آن واحد. نعتبر الحوادث:
\begin{itemize}
\item $R$: «سحب كريتين حمراوين من $U$»،\quad $A$: «سحب كريتين حمراوين من $V$»،
\item $B$: «سحب كريتين خضراوين من $V$»،\quad $C$: «سحب كريتين من لونين مختلفين من $V$».
\end{itemize}
\begin{center}\begin{LTR}\begin{tikzpicture}[bmS,level 1/.style={sibling distance=3.2cm},level 2/.style={sibling distance=1cm},scale=0.9]
\node[ev]{}
 child{node[ev]{$R$} child{node[ev]{$A$} edge from parent node[p]{$\ldots$}} child{node[ev]{$B$} edge from parent node[p]{$\ldots$}} child{node[ev]{$C$} edge from parent node[p]{$\ldots$}} edge from parent node[p]{$\ldots$}}
 child{node[ev]{$\overline{R}$} child{node[ev]{$A$} edge from parent node[p]{$\ldots$}} child{node[ev]{$B$} edge from parent node[p]{$\ldots$}} child{node[ev]{$C$} edge from parent node[p]{$\ldots$}} edge from parent node[p]{$\ldots$}};
\end{tikzpicture}\end{LTR}\end{center}
\QQ{1} أنقل واملأ شجرة الاحتمالات المقابلة.
\QQ{2} تحقق أن $P(A)=\dfrac{7}{30}$ واحسب $P(B)$ و$P(C)$.
\QQ{3} احسب احتمال سحب كريتين حمراوين من $U$ علما أن الكريتين المسحوبتين من $V$ خضراوان.
\QQ{4} $X$ المتغير العشوائي الذي يرفق بكل عملية سحب بالكيفية السابقة عدد الكريات الخضراء المسحوبة من الكيس $V$.
عيّن قانون احتمال المتغير العشوائي $X$ ثم احسب الأمل الرياضياتي $E(21X+1428)$.
\end{tamrin}
\begin{molahada}
في نص الموضوع: «عدد الكريات الخضراء المسحوبة من الكيس»؛ المقصود الكيس $V$ (السحب الأخير)، وهو ما يعطي أملا رياضياتيا صحيحا $E(21X+1428)=1447$.
\end{molahada}

%======================================================================
\begin{tamrin}{بكالوريا 2026 -- الموضوع الثاني -- شعبة رياضيات}
يضم قسم لشعبة الرياضيات $8$ ذكور من بينهم أحمد و$6$ إناث من بينهن فاطمة.

\QQ{1} نختار بطريقة عشوائية $3$ تلاميذ من القسم للمشاركة في أولمبياد الرياضيات ونعتبر الحوادث التالية:
\begin{itemize}
\item $A$: «التلاميذ المختارون من نفس الجنس»،\quad $B$: «من بين التلاميذ المختارين أنثى على الأقل»،
\item $C$: «فاطمة من بين التلاميذ المختارين»،\quad $D$: «أحمد وفاطمة لا يتم اختيارهما معا»،
\item $E$: «فاطمة من بين التلاميذ المختارين علما أن كل التلاميذ المختارين من الإناث».
\end{itemize}
احسب احتمالات الحوادث $A$، $B$، $C$، $D$ و$A\cup B$ وبيّن أن احتمال الحادثة $E$ يساوي $\dfrac12$.

\QQ{2} نفرض أن عدد الإناث هو $n$ حيث $n\geqslant 2$ ونختار عشوائيا ممثلين اثنين للقسم.
$X$ المتغير العشوائي الذي يرفق بكل عملية اختيار لممثلي القسم عدد الإناث المختارات و$E(X)$ أمله الرياضياتي.
\begin{sq}
\qa عيّن قانون احتمال المتغير العشوائي $X$ ثم بيّن أن: $E(X)=\dfrac{2n}{n+8}$.
\qb جد قيمة $n$ التي من أجلها يكون $E(X)$ عددا طبيعيا.
\end{sq}
\end{tamrin}

%======================================================================
\begin{tamrin}{بكالوريا 2025 -- الموضوع الأول -- شعبة رياضيات}
يحتوي صندوق $U_1$ على $12$ كرية منها: $3$ كريات حمراء، $4$ كريات بيضاء و$5$ كريات خضراء.
ويحتوي صندوق $U_2$ على $7$ كريات منها: $4$ كريات حمراء و$3$ كريات بيضاء.
(جميع الكريات متماثلة ولا نفرق بينها باللمس)

نرمي نردا غير مزيف ذا $6$ أوجه مرقمة بـ: $1$، $1$، $1$، $1$، $2$، $2$.
\begin{itemize}
\item إذا ظهر الرقم $1$، نسحب عشوائيا $3$ كريات في آن واحد من $U_1$.
\item وإذا ظهر الرقم $2$، نسحب عشوائيا $3$ كريات على التوالي من $U_2$ دون إرجاع.
\end{itemize}
نعتبر الحوادث: $A$: «ظهور الرقم $1$»، $B$: «الحصول على $3$ كريات من نفس اللون»،
$C$: «الحصول على $3$ كريات مختلفة الألوان مثنى مثنى»، $D$: «الحصول على كريتين بالضبط من نفس اللون».
\begin{center}\begin{LTR}\begin{tikzpicture}[bmS,level 1/.style={sibling distance=2.8cm},level 2/.style={sibling distance=1cm},scale=0.9]
\node[ev]{}
 child{node[ev]{$A$} child{node[ev]{$B$} edge from parent node[p]{$\ldots$}} child{node[ev]{$C$} edge from parent node[p]{$\ldots$}} child{node[ev]{$D$} edge from parent node[p]{$\ldots$}} edge from parent node[p]{$\ldots$}}
 child{node[ev]{$\overline{A}$} child{node[ev]{$B$} edge from parent node[p]{$\ldots$}} child{node[ev]{$D$} edge from parent node[p]{$\ldots$}} edge from parent node[p]{$\ldots$}};
\end{tikzpicture}\end{LTR}\end{center}
\QQ{1} أنقل واملأ شجرة الاحتمالات المقابلة.
\QQ{2}
\begin{sq}
\qa احسب $P(B)$، $P(C)$، $P(D)$، $P(B\cup C)$ احتمالات الحوادث $B$، $C$، $D$، $B\cup C$ على الترتيب.
\qb بيّن أن: $P(A\cap\overline{B})=\dfrac{41}{66}$.
\qc احسب $P_{\overline{B}}(A)$ احتمال سحب $3$ كريات من $U_1$ علما أنها ليست من نفس اللون.
\end{sq}
\QQ{3} $X$ المتغير العشوائي الذي يرفق بكل عملية سحب لـ $3$ كريات كما سبق، عدد الألوان التي تحملها الكريات المسحوبة.
عيّن قانون الاحتمال للمتغير العشوائي $X$ ثم احسب أمله الرياضياتي.
\end{tamrin}

%======================================================================
\begin{tamrin}{بكالوريا 2024 -- الموضوع الأول -- شعبة رياضيات}
يحتوي صندوق $U_1$ على $5$ بطاقات متماثلة مرقمة بـ: $1$، $1$، $2$، $3$، $3$ ويحتوي صندوق $U_2$ على $6$ كريات موزعة كما يلي: $4$ كريات حمراء و$2$ خضراوان. (لا نفرق بين البطاقات ولا بين الكريات باللمس)

نسحب عشوائيا بطاقة واحدة من الصندوق $U_1$.
\begin{itemize}
\item إذا تحصلنا على الرقم $1$ نسحب عشوائيا من $U_2$ كرية واحدة.
\item إذا تحصلنا على الرقم $2$ نسحب عشوائيا من $U_2$ كريتين في آن واحد.
\item إذا تحصلنا على الرقم $3$ نسحب عشوائيا من $U_2$ ثلاث كريات في آن واحد.
\end{itemize}
نعتبر الحوادث الآتية: $C_i$: «البطاقة المتحصل عليها تحمل الرقم $i$» حيث $i\in\{1;2;3\}$،
$A$: «الحصول على كريات حمراء فقط»، $B$: «الحصول على كريات خضراء فقط»، $D$: «الحصول على كريات ليست كلها من نفس اللون».
\begin{center}\begin{LTR}\begin{tikzpicture}[bmS,level 1/.style={sibling distance=2.8cm},level 2/.style={sibling distance=1cm},scale=0.9]
\node[ev]{}
 child{node[ev]{$C_1$} child{node[ev]{$A$} edge from parent node[p]{$\ldots$}} child{node[ev]{$B$} edge from parent node[p]{$\ldots$}} edge from parent node[p]{$\ldots$}}
 child{node[ev]{$C_2$} child{node[ev]{$A$} edge from parent node[p]{$\ldots$}} child{node[ev]{$B$} edge from parent node[p]{$\ldots$}} child{node[ev]{$D$} edge from parent node[p]{$\ldots$}} edge from parent node[p]{$\ldots$}}
 child{node[ev]{$C_3$} child{node[ev]{$A$} edge from parent node[p]{$\ldots$}} child{node[ev]{$D$} edge from parent node[p]{$\ldots$}} edge from parent node[p]{$\ldots$}};
\end{tikzpicture}\end{LTR}\end{center}
\QQ{1}
\begin{sq}
\qa بيّن أن: $P_{C_2}(B)=\dfrac{1}{15}$ و$P_{C_3}(D)=\dfrac45$.
\qb أنقل واملأ شجرة الاحتمالات المقابلة.
\qc احسب $P(A)$، $P(B)$ و$P(D)$.
\end{sq}
\QQ{2} احسب احتمال أن تكون البطاقة المتحصل عليها تحمل الرقم $3$ علما أن الكريات المسحوبة حمراء.
\QQ{3} $X$ المتغير العشوائي الذي يرفق بكل عملية سحب عدد الألوان المتحصل عليها.
عيّن قانون الاحتمال للمتغير العشوائي $X$ ثم احسب الأمل الرياضياتي $E(75X+1917)$.
\QQ{4} إذا كان عدد الكريات الحمراء في الصندوق $U_2$ هو $n+4$ حيث $n$ عدد طبيعي، جد قيمة $n$ التي من أجلها يكون $P_{C_3}(A)=\dfrac{7}{15}$.
\end{tamrin}

%======================================================================
\begin{tamrin}{بكالوريا 2024 -- الموضوع الثاني -- شعبة رياضيات}
يحتوي صندوق $U_1$ على $7$ كريات منها: $3$ كريات بيضاء و$4$ كريات حمراء، ويحتوي صندوق $U_2$ على $7$ كريات منها: كريتان بيضاوان و$5$ كريات حمراء. (جميع الكريات متماثلة ولا نفرق بينها باللمس)

نلقي نردا متوازنا أوجهه مرقمة من $1$ إلى $6$ ونسجل الرقم الظاهر على الوجه العلوي.
\begin{itemize}
\item إذا ظهر رقم مضاعف للعدد $3$، نسحب عشوائيا من الصندوق $U_1$ كريتين على التوالي دون إرجاع.
\item في الحالات الأخرى، نسحب عشوائيا من الصندوق $U_2$ كريتين على التوالي دون إرجاع.
\end{itemize}
نعتبر الحوادث الآتية: $M$: «ظهور رقم مضاعف للعدد $3$»، $A$: «الحصول على كريتين بيضاوين»،
$B$: «الحصول على كريتين حمراوين»، $C$: «الحصول على كريتين من لونين مختلفين».
\begin{center}\begin{LTR}\begin{tikzpicture}[bmS,level 1/.style={sibling distance=3.2cm},level 2/.style={sibling distance=1cm},scale=0.9]
\node[ev]{}
 child{node[ev]{$M$} child{node[ev]{$A$} edge from parent node[p]{$\ldots$}} child{node[ev]{$B$} edge from parent node[p]{$\ldots$}} child{node[ev]{$C$} edge from parent node[p]{$\ldots$}} edge from parent node[p]{$\ldots$}}
 child{node[ev]{$\overline{M}$} child{node[ev]{$A$} edge from parent node[p]{$\ldots$}} child{node[ev]{$B$} edge from parent node[p]{$\ldots$}} child{node[ev]{$C$} edge from parent node[p]{$\ldots$}} edge from parent node[p]{$\ldots$}};
\end{tikzpicture}\end{LTR}\end{center}
\QQ{1} أنقل واملأ شجرة الاحتمالات المقابلة.
\QQ{2} نعتبر الحادثتين: $G$: «الحصول على كريتين من نفس اللون» و$H$: «الحصول على كرية حمراء على الأقل».
بيّن أن $P(G)=\dfrac{31}{63}$ ثم احسب $P(H)$.
\QQ{3} احسب $P_G(M)$ احتمال ظهور رقم مضاعف للعدد $3$ علما أن الكريتين المسحوبتين من نفس اللون.
\QQ{4} $X$ المتغير العشوائي الذي يرفق بكل عملية سحب لكريتين عدد الألوان المتحصل عليها.
عيّن قانون الاحتمال للمتغير العشوائي $X$ ثم احسب الأمل الرياضياتي $E(63X+1350)$.
\end{tamrin}

%======================================================================
\begin{tamrin}{بكالوريا 2023 -- الموضوع الأول -- شعبة رياضيات}
يحتوي كيس على $10$ كريات متماثلة ولا نفرق بينها باللمس، منها كريتان حمراوان مرقمتان بـ: $-3$، $2$ وخمس كريات بيضاء مرقمة بـ: $-2$، $0$، $1$، $1$، $2$ وثلاث كريات خضراء مرقمة بـ: $0$، $1$، $2$.
نسحب عشوائيا وفي آن واحد ثلاث كريات من الكيس ونعتبر الأحداث $A$، $B$ و$C$ الآتية:
$A$: «الحصول على $3$ كريات من نفس اللون»، $B$: «الحصول على الألوان الثلاثة»، $C$: «الحصول على $3$ كريات مجموع أرقامها معدوم».

\QQ{1}
\begin{sq}
\qa احسب $P(A)$ و$P(B)$ ثم بيّن أن: $P(C)=\dfrac{3}{20}$.
\qb احسب $P(A\cap C)$ ثم استنتج $P_C(A)$.
\end{sq}
\QQ{2} نعتبر المتغير العشوائي $X$ الذي يرفق بكل عملية سحب لثلاث كريات عدد الألوان المتحصل عليها.
عيّن قانون الاحتمال للمتغير العشوائي $X$ ثم احسب أمله الرياضياتي $E(X)$.
\QQ{3} نسحب الآن عشوائيا من الكيس ثلاث كريات على التوالي وبإرجاع.
احسب احتمال الحصول على ثلاث كريات جداء أرقامها معدوم.
\end{tamrin}
\begin{molahada}
في النص: «احسب $P(A)$ و$P(C)$ ثم بيّن أن $P(C)=\frac{3}{20}$»؛ الصواب: احسب $P(A)$ و$P(B)$.
\end{molahada}

%======================================================================
\begin{tamrin}{بكالوريا 2023 -- الموضوع الثاني -- شعبة رياضيات}
يحتوي صندوق $U$ على كريتين حمراوين وكريتين خضراوين ويحتوي صندوق $V$ على كريتين حمراوين وثلاث كريات خضراء. (كل الكريات متماثلة لا نفرق بينها عند اللمس)

نسحب عشوائيا كريتين في آن واحد من أحد الصندوقين بالكيفية الآتية:
نقوم بسحب بطاقة واحدة عشوائيا من كيس به $10$ بطاقات متماثلة ومرقمة من $1$ إلى $10$.
إذا تحصلنا على عدد أولي نسحب الكريتين من $U$ وفي باقي الحالات نسحب الكريتين من $V$.

\QQ{1} نعتبر الأحداث $A$، $B$ و$C$ الآتية: $A$: «سحب كريتين حمراوين»، $B$: «سحب كريتين خضراوين» و$C$: «سحب كريتين من لونين مختلفين».
\begin{sq}
\qa أنجز شجرة الاحتمالات التي تنمذج هذه التجربة.
\qb بيّن أن $P(A)=\dfrac{19}{150}$ و$P(B)=\dfrac{37}{150}$ ثم استنتج $P(C)$.
\end{sq}
\QQ{2} $X$ المتغير العشوائي الذي يرفق بكل عملية سحب لكريتين عدد الكريات الحمراء المتحصل عليها.
\begin{sq}
\qa عيّن قانون الاحتمال للمتغير العشوائي $X$ ثم احسب أمله الرياضياتي $E(X)$.
\qb احسب احتمال الحدث: «$\ln X\leqslant 1$».
\end{sq}
\end{tamrin}
\begin{molahada}
في النص: «بيّن أن $P(B)=\frac{19}{150}$ و$P(B)=\frac{37}{150}$»؛ الصواب: $P(A)=\frac{19}{150}$ و$P(B)=\frac{37}{150}$.
\end{molahada}

%======================================================================
\begin{tamrin}{بكالوريا 2021 -- الموضوع الأول -- شعبة رياضيات}
كيس به $12$ كرية متماثلة لا نفرق بينها باللمس.
كل من الكريات الاثنتي عشرة تحمل رقما من بين الأعداد التالية: $1$، $2$، $3$ و$4$.
نسحب عشوائيا كرية واحدة من الكيس.
نرمز بـ $P_i$ إلى احتمال سحب كرية رقمها $i$، حيث: $P_1=\dfrac13$، $P_2=\dfrac16$، $P_3=\dfrac14$ و$P_4=\dfrac14$.

\QQ{1} وزّع الكريات الاثنتي عشرة حسب الأرقام $1$، $2$، $3$ و$4$.
\QQ{2} احسب احتمال كل من الأحداث $A$، $B$ و$C$ الآتية:
$A$: «سحب كرية تحمل رقما فرديا»، $B$: «سحب كرية تحمل رقما من أرقام نظام التعداد ذي الأساس $4$»، $C$: «سحب كرية رقمها حل للمعادلة $x^2=2^x$».
\QQ{3} المتغير العشوائي $X$ يرفق بكل سحب لكرية الرقم الذي تحمله.
عيّن مجموعة قيم المتغير العشوائي $X$ ثم احسب $E(X)$ أمله الرياضياتي.
\end{tamrin}

%======================================================================
\begin{tamrin}{بكالوريا 2021 -- الموضوع الثاني -- شعبة رياضيات}
يراد عشوائيا تشكيل لجنة تضم رئيسا ونائبا له من بين ثلاثة رجال $H_1$، $H_2$ و$H_3$ وأربع نساء $F_1$، $F_2$، $F_3$ و$F_4$.

\QQ{1} بيّن أن عدد اللجان التي يمكن تشكيلها هو: $42$.
\QQ{2} نعتبر الأحداث الآتية:
$A$: «اللجنة من نفس الجنس»، $B$: «اللجنة من جنسين مختلفين»، $C$: «$H_1$ هو الرئيس»، $E$: «اللجنة لا تضم كلا من $H_1$ و$F_1$».
\begin{sq}
\qa احسب $P(A)$ احتمال الحدث $A$ ثم استنتج $P(B)$.
\qb احسب $P(C)$ و$P(E)$.
\end{sq}
\QQ{3} المتغير العشوائي $X$ يرفق بكل لجنة عدد الرجال فيها.
عيّن قانون احتمال $X$ ثم احسب $E(X)$ أمله الرياضياتي.
\end{tamrin}

%======================================================================
\begin{tamrin}{بكالوريا 2020 -- الموضوع الأول -- شعبة رياضيات}
صندوق به $5$ كريات بيضاء و$3$ كريات حمراء. (كل الكريات متماثلة لا نفرق بينها باللمس)
نسحب من الصندوق كرية واحدة بحيث: إذا ظهرت كرية حمراء نعيدها إلى الصندوق ونضيف له كرية بيضاء، وإذا ظهرت كرية بيضاء نعيدها إلى الصندوق ونضيف له كرية حمراء، ونكرر العملية مرة ثانية.
\begin{center}\begin{LTR}\begin{tikzpicture}[bmS,level 1/.style={sibling distance=2.8cm},level 2/.style={sibling distance=1cm},scale=0.9]
\node[ev]{}
 child{node[ev]{$R$} child{node[ev]{$R$} edge from parent node[p]{$\ldots$}} child{node[ev]{$B$} edge from parent node[p]{$\ldots$}} edge from parent node[p]{$\ldots$}}
 child{node[ev]{$B$} child{node[ev]{$R$} edge from parent node[p]{$\frac49$}} child{node[ev]{$B$} edge from parent node[p]{$\ldots$}} edge from parent node[p]{$\ldots$}};
\end{tikzpicture}\end{LTR}\end{center}
\QQ{1} أنقل شجرة الاحتمالات المقابلة التي تنمذج هذه التجربة ثم أكملها.
\QQ{2} بيّن أن احتمال أن يوجد في الصندوق $7$ كريات بيضاء هو $\dfrac18$.
\QQ{3} احسب احتمال أن يوجد بالصندوق $4$ كريات حمراء على الأقل.
\QQ{4} ليكن $X$ المتغير العشوائي الذي يأخذ كقيمة عدد الكريات البيضاء الموجودة في الصندوق بعد العملية الثانية.
\begin{sq}
\qa برّر أن قيم المتغير العشوائي هي: $5$، $6$ و$7$.
\qb عرّف قانون الاحتمال للمتغير العشوائي $X$، ثم احسب $E(X)$ أمله الرياضياتي.
\end{sq}
\end{tamrin}

%======================================================================
\begin{tamrin}{بكالوريا 2020 -- الموضوع الثاني -- شعبة رياضيات}
يحتوي صندوق على كريات متماثلة منها: $n$ كرية بيضاء تحمل العدد $\pi$ ($n$ عدد طبيعي و$n\geqslant2$) و$4$ كريات حمراء تحمل الأعداد $\frac{\pi}{2}$، $\frac{\pi}{2}$، $\frac{\pi}{3}$ و$\pi$ وكريتان خضراوان تحملان العددين $\frac{\pi}{2}$ و$\frac{\pi}{3}$.
نسحب عشوائيا كريتين في آن واحد من هذا الصندوق.

\QQ{1}
\begin{sq}
\qa احسب احتمال كل من الحدثين $A$ و$B$ حيث: $A$: «سحب كريتين من نفس اللون» و$B$: «سحب كريتين تحملان نفس العدد علما أنهما من نفس اللون».
\qb عيّن العدد الطبيعي $n$ حتى يكون: $P(A)=\dfrac{17}{55}$.
\end{sq}
\QQ{2} نفرض فيما يلي: $n=5$ ونسمي $\alpha$ و$\beta$ العددين الظاهرين على الكريتين المسحوبتين.
نعتبر $X$ المتغير العشوائي الذي يرفق بكل نتيجة سحب العدد: $\cos(\alpha)\cos(\beta)$.
\begin{sq}
\qa برّر أن قيم المتغير العشوائي $X$ هي: $-\frac12$، $0$، $\frac14$ و$1$.
\qb بيّن أن: $P(X=0)=\dfrac{27}{55}$.
\qc عرّف قانون الاحتمال للمتغير العشوائي $X$، ثم احسب أمله الرياضياتي $E(X)$.
\end{sq}
\end{tamrin}

%======================================================================
\begin{tamrin}{بكالوريا 2019 -- الموضوع الثاني -- شعبة رياضيات}
صندوقان غير شفافين $U_1$ و$U_2$، يحتوي الصندوق $U_1$ على $4$ كريات حمراء وثلاث كريات سوداء ويحتوي الصندوق $U_2$ على $3$ كريات حمراء وكريتين سوداوين. (كل الكريات متشابهة لا نفرق بينها عند اللمس)

نرمي نردا غير مزيف ذا ستة أوجه مرقمة من $1$ إلى $6$.
إذا ظهر الرقمان $2$ أو $4$ نسحب عشوائيا كريتين في آن واحد من الصندوق $U_1$، وفي باقي الحالات نسحب عشوائيا كريتين في آن واحد من الصندوق $U_2$.
نعتبر الأحداث $A$، $B$ و$C$ المعرفة بـ: $A$: «سحب كريتين حمراوين»، $B$: «سحب كريتين سوداوين»، $C$: «سحب كريتين من لونين مختلفين».
\begin{center}\begin{LTR}\begin{tikzpicture}[bmS,level 1/.style={sibling distance=3.2cm},level 2/.style={sibling distance=1cm},scale=0.9]
\node[ev]{}
 child{node[ev]{$U_1$} child{node[ev]{$A$} edge from parent node[p]{$\frac27$}} child{node[ev]{$B$} edge from parent node[p]{$\frac17$}} child{node[ev]{$C$} edge from parent node[p]{$\ldots$}} edge from parent node[p]{$\frac13$}}
 child{node[ev]{$U_2$} child{node[ev]{$A$} edge from parent node[p]{$\frac3{10}$}} child{node[ev]{$B$} edge from parent node[p]{$\ldots$}} child{node[ev]{$C$} edge from parent node[p]{$\ldots$}} edge from parent node[p]{$\ldots$}};
\end{tikzpicture}\end{LTR}\end{center}
\QQ{1} أنقل وأكمل شجرة الاحتمالات.
\QQ{2} احسب احتمالات الأحداث $A$، $B$ و$C$.

نعتبر المتغير العشوائي $X$ الذي يرفق بكل عملية سحب عدد الكريات الحمراء المسحوبة.
\QQ{3}
\begin{sq}
\qa عيّن قيم المتغير العشوائي $X$.
\qb عرّف قانون الاحتمال للمتغير العشوائي $X$.
\end{sq}
\QQ{4} احسب الأمل الرياضياتي $E(X)$.
\end{tamrin}

%======================================================================
\begin{tamrin}{بكالوريا 2018 -- الموضوع الثاني -- شعبة رياضيات}
كيس يحوي $9$ كريات لا نفرق بينها باللمس موزعة كما يلي:
خمس كريات حمراء مرقمة بـ: $1$، $1$، $2$، $2$، $2$ وثلاث كريات خضراء مرقمة بـ: $-3$، $2$، $3$ وكرية بيضاء مرقمة بـ: $-1$.
نسحب عشوائيا $4$ كريات في آن واحد.

\QQ{1} احسب احتمال الحوادث التالية:
$A$: «الحصول على أربع كريات من نفس اللون»، $B$: «الحصول على كرية بيضاء على الأكثر»، $C$: «الحصول على أربع كريات مجموع أرقامها معدوم».
\QQ{2} ليكن $X$ المتغير العشوائي الذي يرفق بكل نتيجة سحب عدد الكريات الخضراء المتبقية في الكيس.
\begin{sq}
\qa عيّن قيم المتغير العشوائي $X$ ثم عرّف قانون احتماله.
\qb احسب الأمل الرياضياتي $E(X)$ للمتغير العشوائي $X$.
\qc احسب احتمال الحادثة: «$X^2-X>0$».
\end{sq}
\end{tamrin}

%======================================================================
\begin{tamrin}{بكالوريا 2009 -- الموضوع الأول -- شعبة رياضيات}
كيس به $10$ كريات متماثلة لا نميز بينها عند اللمس منها $4$ بيضاء و$6$ حمراء.

\QQ{1} نسحب عشوائيا من الكيس $3$ كريات في آن واحد.
\begin{sq}
\qa احسب احتمال الحصول على $3$ كريات بيضاء.
\qb احسب احتمال الحصول على الأقل على كرية حمراء.
\end{sq}
\QQ{2} ليكن $X$ المتغير العشوائي الذي يرفق بكل نتيجة سحب عدد الكريات البيضاء المسحوبة.
عرّف قانون الاحتمال للمتغير العشوائي $X$ واحسب أمله الرياضياتي.
\end{tamrin}

\end{document}
